\documentclass[12pt]{article}

\input{../../_assets/preambulo.tex}


\begin{document}

    % 1. Foto de fondo
    % 2. Título
    % 3. Encabezado Izquierdo
    % 4. Color de fondo
    % 5. Coord x del titulo
    % 6. Coord y del titulo
    % 7. Fecha

    
    \input{../../_assets/portada}
    \portadaExamen{ffccA4.jpg}{Ecuaciones\\Diferenciales II\\Examen VI}{Ecuaciones Diferenciales II. Examen VI}{MidnightBlue}{-8}{28}{2026}{}

    \begin{description}
        \item[Asignatura] Ecuaciones Diferenciales II.
        \item[Curso Académico] 2021/2022.
        \item[Grado] Doble Grado en Ingeniería Informática y Matemáticas
        \item[Grupo] Único
        \item[Profesor] Rafael Ortega Ríos
        \item[Descripción] Parcial 1.
        \item[Fecha] $4$ de Abril de $2022$.
        % \item[Duración] ---.
    
    \end{description}
    \newpage


    % ------------------------------------
    
    \begin{ejercicio}
        Dada una norma $\|\cdot\|$ en $\bb{R}^d$, construye una función continua del tipo $\psi : \bb{R}^d \to [0, 1]$ que cumpla
        \begin{equation*}
            \psi(x) = \begin{cases}
                0, & \|x\| \leq 3,\\
                1, & \|x\| \geq 7.
            \end{cases}
        \end{equation*}
    \end{ejercicio}

    \begin{ejercicio}~
        \begin{enumerate}[label=\alph*)]
            \item Prueba la desigualdad
            \[
            \sqrt{a+b} \leq \sqrt{a} + \sqrt{b}, \quad \forall a, b \in [0, \infty[.
            \]
            \item Demuestra que la sucesión de funciones $\{f_n\}_{n \geq 1}$,
            \[
            f_n(t) = \sqrt{t + \frac{1}{n}}, \quad t \in [0, 1]
            \]
            converge uniformemente en $[0, 1]$.
            \item ¿Se puede afirmar que también la sucesión de derivadas $\{f'_n\}_{n \geq 1}$ es uniformemente convergente en $[0, 1]$?
        \end{enumerate}
    \end{ejercicio}

    \begin{ejercicio}
        Se considera el sistema de ecuaciones integrales
        \[
        x_1(t) = -\int_{0}^{t} x_2(s)ds, \quad x_2(t) = \int_{\pi}^{t} x_1(s)ds.
        \]
        ¿Existe solución continua y definida en $]-\infty, \infty[$? ¿Es única?
    \end{ejercicio}

    \begin{ejercicio}
        Se considera el problema de valores iniciales
        \[
        x' = \frac{\sqrt{3 - t^2}}{1 + x^2}, \quad x(0) = 0.
        \]
        ¿Existe solución? ¿Es única? En caso afirmativo determina el intervalo maximal.
    \end{ejercicio}

    \begin{ejercicio}
        La función $F : [-1, 1] \times \bb{R} \to \bb{R}$, $(t, x) \mapsto F(t, x)$ es continua y cumple la condición
        \[
        |F(t, x_1) - F(t, x_2)| \leq L|x_1 - x_2|, \quad \forall(t, x) \in [-1, 1] \times \bb{R},
        \]
        con $L < \nicefrac{1}{2}$. Demuestra que existe a lo sumo una función continua $x : [-1, 1] \to \bb{R}$ que cumple la ecuación integral
        \[
        x(t) = \int_{-t}^{t} F(s, x(s))ds.
        \]
    \end{ejercicio}

    \newpage
    \setcounter{ejercicio}{0} % Reiniciar contador de ejercicios
    % \noindent
    % \textbf{Solución.}

\end{document}